\begin{tcolorbox}[colback=referenceBlueBg,colframe=genderMasculine,title={B1 coverage note}]
This reference contains 635 complete verb tables. A further 334 separable
compounds are listed compactly under their stems, giving 969 indexed verb
headwords. The selection focuses on practical vocabulary through CEFR B1,
with selected related verbs for broader everyday reading; CEFR
does not prescribe one official exhaustive verb list. Each full table includes
all applicable present and simple-past forms, four imperative forms where they
exist, and the perfect, future and pluperfect constructions.
\end{tcolorbox}

\begin{tcolorbox}[colback=white,colframe=genderMasculine,title={Related verbs to compare},fonttitle=\bfseries]
These groups connect entries in the alphabetical and reflexive chapters.
Related words can differ in meaning, separability, case government and auxiliary.

\smallskip
\begin{tabularx}{\linewidth}{@{}lX@{}}
Tension and stretching & \textit{spannen, anspannen, entspannen, dehnen, strecken}\\
Writing and prescribing & \textit{schreiben, aufschreiben, verschreiben, vorschreiben, verordnen}\\
Greetings and farewells & \textit{begrüßen, verabschieden, sich verabschieden}\\
Driving and pursuing & \textit{treiben, antreiben, betreiben, vertreiben, übertreiben}\\
Handing over & \textit{reichen, überreichen, übergeben, weitergeben, aushändigen}\\
Restraint and restriction & \textit{binden, fesseln, festhalten, hindern, einschränken, beschränken}\\
Catching and detention & \textit{fangen, erwischen, festnehmen, verhaften, einsperren, inhaftieren}\\
Freedom and release & \textit{befreien, freilassen, entlassen}\\
\end{tabularx}

\smallskip
\textbf{Usage contrasts.}
\textit{Sport treiben} means to do sport; \textit{ein Geschäft betreiben}
means to run a business. \textit{Verschreiben} usually concerns medicine,
whereas \textit{vorschreiben} can set a rule:
\textit{Der Arzt verschreibt mir Tabletten. Das Gesetz schreibt es vor.}

\textit{Jemanden verabschieden} means to see someone off;
\textit{sich von jemandem verabschieden} means to say goodbye to someone.
\textit{Sich übergeben} means to vomit, rather than to hand something over.
Compare the inseparable \textit{übergeben/überreichen}
(\textit{hat übergeben/überreicht}) with the separable \textit{weitergeben}
(\textit{hat weitergegeben}).
\end{tcolorbox}
